
% Abstract File
% Do not remove centering environment below
\begin{center}

\end{center}

\begin{center}
{{\bf\fontsize{14pt}{14.5pt}\selectfont \uppercase{ABSTRACT}}}
\end{center}

\doublespacing
\addcontentsline{toc}{chapter}{Abstract}



% \begin{center}
% 	\begin{doublespace}
% {\fontsize{14pt}{14.5pt}\selectfont {Title of Dissertation/Thesis}}\\
% 		{\fontsize{14pt}{14.5pt}\selectfont {Student Full Name}}\\
% %		{\fontsize{12pt}{12.5pt}\selectfont {Ph.D. Year}}\\
% 	\end{doublespace}
% \end{center}


% \textcolor{blue}{[wolę napisać w miarę na końcu - teraz jest to tekst z wniosku.] The research will focus on the comparative evaluation of various machine learning and deep learning models for classifying human-written versus machine-generated text. The study will explore a range of methods, including traditional machine learning techniques based on feature engineering such as Logistic Regression (LR) and Support Vector Machines (SVM), as well as more advanced deep learning architectures like convolutional neural networks (CNNs), recurrent neural networks (RNNs), and transformers such as RoBERTa. By training these models on specific datasets and analyzing their performance metrics, the research aims to highlight the strengths, weaknesses, and limitations of each approach. Explainable AI (XAI) techniques will be used to interpret and explain the models' decision-making processes. Special attention will be given to how different models perform with different types of data, identifying the most suitable approaches for varying classification tasks. A prototype application will also be developed, enabling users to assess whether a given text is AI-generated or human-written based on the trained models.}


This thesis develops and evaluates a hybrid multi-agent system for algorithmic trading, in which each agent relies on a distinct artificial-intelligence technique. It draws on OHLCV price data, technical indicators, microstructure and market-structure features, and external sentiment and cross-asset signals, and is trained and evaluated on hourly Bitcoin (BTC/USDT) data under a strict, leakage-controlled walk-forward protocol.

\vspace{3mm}

Four learned models from different paradigms were built --- a gradient-boosting model (LightGBM), a temporal convolutional network (TCN), a patch-based Transformer (PatchTST), and a selective state-space model (Mamba) --- alongside several rule-based agents. The final fund is a \textit{predeclared roster} of five autonomous agents, curated from this pool for paradigm diversity and acceptable validation-window behaviour: three learned models (LightGBM, TCN, Mamba) and two rule-based agents (trend following and cross-asset dominance rotation). PatchTST is built and analysed but \textit{excluded}: its directional accuracy is no better than chance out-of-sample --- a result that persists under walk-forward retraining and instance normalisation --- so it contributes no forecasting skill. Each roster agent produces its own risk-managed stream of returns, combined by a leak-free, capped inverse-volatility allocator. Other approaches explored along the way --- deep reinforcement learning, evolutionary rule generation, a volatility-breakout rule, a contrarian sentiment rule, and recurrent networks --- did not make the final fund and are reported as honest negatives.

\vspace{3mm}

Over a two-year out-of-sample window spanning three market regimes, the diversified fund achieved the strongest \textit{risk-adjusted} performance among the strategies it comprises: a $+53.7\%$ return at a Sharpe ratio of $1.50$ and a maximum drawdown of $-9.2\%$, strongly positive in the bull and bear regimes and flat in chop, ahead of Bitcoin buy-and-hold (Sharpe $0.09$) and the S\&P~500 (Sharpe $1.16$). The result is robust but modest: a block bootstrap places the fund's Sharpe firmly above zero, yet a null of random weightings of the same five agents shows the capped allocator barely edges an equal-weight blend (Sharpe $1.48$), so the value lies in \textit{holding} a diverse roster rather than in the allocation rule. A learned adaptive coordinator failed outright to beat the simple risk-parity rule, and the highest-returning single agent did not survive a random-bracket test for genuine skill. The central conclusion is that combining heterogeneous AI methods \textit{can} yield more stable, regime-balanced performance than any single approach, with the value lying in disciplined diversification across genuinely diverse agents rather than in a sophisticated controller. Live deployment --- with its real-world execution costs --- is the natural next step.
